% !TEX root =./ECE343_Practice_main.tex
%
% L25 practice set (Boundary Conditions and Capacitance): the parallel-plate
% capacitor, the coaxial C' derived and applied to RG-59, and a capacitor with
% two dielectric layers.

\fullwidth{\textbf{Learning Objective 5.4: } I can apply the electric field
boundary conditions to calculate electric fields on both sides of a material
boundary.
\\[4pt]\normalfont\small Show your work. Write a \textbf{Documentation} line at the
end (\textbf{None} if you did not collaborate). Ignore fringing in every
parallel-plate problem.}

\begin{questions}

%%%%%%%%%%%%%%%%%%%%%%%%%%%%%%%%%%%%%%%%%
\question \textbf{Parallel plates.} Two plates of area $A$ are a distance $d$
apart, with a dielectric of permittivity $\epsilon$ between them. The plate at
$z = d$ carries $+\rho_s$ and the plate at $z = 0$ carries $-\rho_s$.
\begin{parts}

	\part Use the conductor boundary condition to find $\mathbf{E}$ between the
	plates, then integrate it to derive $C = \epsilon A/d$.
	\begin{solution}[2in]
	At the top plate the outward normal into the gap is $-\hat{\mathbf{z}}$:
	\eq{
		\mathbf{D}\mycdot\lp-\hat{\mathbf{z}}\rp &= \rho_s \implies \mathbf{E} = -\hat{\mathbf{z}}\dfrac{\rho_s}{\epsilon} \\
		V &= -\int_0^d \lp-\hat{\mathbf{z}}\dfrac{\rho_s}{\epsilon}\rp\mycdot\hat{\mathbf{z}}\,dz = \dfrac{\rho_s d}{\epsilon} \\
		C &= \dfrac{Q}{V} = \dfrac{\rho_s A}{\rho_s d/\epsilon} = \dfrac{\epsilon A}{d}
	}
	\end{solution}

	\begin{minipage}[t]{0.58\linewidth}
		\part The plates measure $2\ \text{cm} \times 3\ \text{cm}$ and are
		$0.5$ mm apart, filled with a film of $\epsilon_r = 3.2$. Find $C$.
		\begin{solution}[0.9in]
		\eq{ C = \dfrac{3.2\times8.854\times10^{-12}\times6\times10^{-4}}{5\times10^{-4}} = 34.0\ \text{pF} }
		\end{solution}
	\end{minipage}\hfill%
	\begin{minipage}[t]{0.37\linewidth}
		\raggedleft \vspace{12pt}
		$C$ = \ansbox[1.3in]{$34.0$ pF}
	\end{minipage}

	\begin{minipage}[t]{0.58\linewidth}
		\part The capacitor is charged to $100$ V. Find $Q$, $\rho_s$, and
		$|\mathbf{E}|$, and check that $|\mathbf{E}| = \rho_s/\epsilon$.
		\begin{solution}[1.6in]
		\eq{
			Q &= CV = 3.40\ \text{nC} \\
			\rho_s &= \dfrac{Q}{A} = 5.67\ \mu\text{C/m}^2 \\
			|\mathbf{E}| &= \dfrac{V}{d} = 200\ \text{kV/m} \\
			\dfrac{\rho_s}{\epsilon} &= \dfrac{5.67\times10^{-6}}{2.83\times10^{-11}} = 200\ \text{kV/m}
		}
		\end{solution}
	\end{minipage}\hfill%
	\begin{minipage}[t]{0.37\linewidth}
		\raggedleft \vspace{12pt}
		$Q$ = \ansbox[1.3in]{$3.40$ nC}

		\vspace{4pt}
		$\rho_s$ = \ansbox[1.3in]{$5.67\ \mu$C/m$^2$}

		\vspace{4pt}
		$|\mathbf{E}|$ = \ansbox[1.3in]{$200$ kV/m}
	\end{minipage}

\end{parts}

\newpage
%%%%%%%%%%%%%%%%%%%%%%%%%%%%%%%%%%%%%%%%%
\question \textbf{Coaxial capacitance.} A coaxial line has inner radius $a$,
shield radius $b$, and a dielectric of permittivity $\epsilon$.
\begin{parts}

	\part Place $+\rho_\ell$ per meter on the inner conductor. Use Gauss's law on a
	cylinder of radius $\rho$ and length $\ell$ to find $\mathbf{E}$ for $a < \rho < b$,
	then derive $C^{\prime} = 2\pi\epsilon/\ln(b/a)$.
	\begin{solution}[2.3in]
	The field is radial and depends only on $\rho$; the end caps carry no flux.
	\eq{
		D_\rho\lp 2\pi\rho\ell\rp &= \rho_\ell\,\ell \implies \mathbf{E} = \hat{\boldsymbol{\rho}}\dfrac{\rho_\ell}{2\pi\epsilon\rho} \\
		V &= -\int_b^a \dfrac{\rho_\ell}{2\pi\epsilon\rho}\,d\rho = \dfrac{\rho_\ell}{2\pi\epsilon}\ln\dfrac{b}{a} \\
		C^{\prime} &= \dfrac{Q/\ell}{V} = \dfrac{\rho_\ell}{\frac{\rho_\ell}{2\pi\epsilon}\ln(b/a)} = \dfrac{2\pi\epsilon}{\ln(b/a)}
	}
	\end{solution}

	\begin{minipage}[t]{0.58\linewidth}
		\part RG-59 has $a = 0.29$ mm, $b = 1.85$ mm, and solid polyethylene,
		$\epsilon_r = 2.25$. Find $C^{\prime}$ and the capacitance of a $15$ m run.
		\begin{solution}[1.1in]
		\eq{
			C^{\prime} &= \dfrac{2\pi\times2.25\times8.854\times10^{-12}}{\ln(1.85/0.29)} \\
			&= \dfrac{1.252\times10^{-10}}{1.853} = 67.5\ \text{pF/m} \\
			C &= 67.5\times15 = 1.01\ \text{nF}
		}
		\end{solution}
	\end{minipage}\hfill%
	\begin{minipage}[t]{0.37\linewidth}
		\raggedleft \vspace{12pt}
		$C^{\prime}$ = \ansbox[1.3in]{$67.5$ pF/m}

		\vspace{4pt}
		$C$ = \ansbox[1.3in]{$1.01$ nF}
	\end{minipage}

	\begin{minipage}[t]{0.58\linewidth}
		\part Module 3 shows that $L^{\prime} = (\mu_0/2\pi)\ln(b/a)$. Use it to find
		$Z_0$ for RG-59, and explain why the cable is sold as $75\ \Omega$ line.
		\begin{solution}[1.2in]
		\eq{
			L^{\prime} &= 2\times10^{-7}\times1.853 = 371\ \text{nH/m} \\
			Z_0 &= \sqrt{\dfrac{L^{\prime}}{C^{\prime}}} = \sqrt{\dfrac{3.71\times10^{-7}}{6.75\times10^{-11}}} = 74.1\ \ohms
		}
		That is within 2\% of $75\ \ohms$.
		\end{solution}
	\end{minipage}\hfill%
	\begin{minipage}[t]{0.37\linewidth}
		\raggedleft \vspace{12pt}
		$Z_0$ = \ansbox[1.3in]{$74.1\ \ohms$}
	\end{minipage}

\end{parts}

\newpage
%%%%%%%%%%%%%%%%%%%%%%%%%%%%%%%%%%%%%%%%%
\question \textbf{Two layers.} Plates of area $0.01\ \text{m}^2$ hold two
stacked layers: $1$ mm with $\epsilon_{r1} = 2$ and $1$ mm with
$\epsilon_{r2} = 8$. The capacitor is charged to $500$ V.
\begin{parts}

	\part Explain why $\mathbf{D}$ is the same in both layers.
	\begin{solution}[0.8in]
	The field is normal to the interface, and the interface carries no free charge,
	so the normal boundary condition makes $D_{n1} = D_{n2}$.
	\end{solution}

	\begin{minipage}[t]{0.58\linewidth}
		\part Find $C$, and check it as two capacitors in series.
		\begin{solution}[1.9in]
		\eq{
			C &= \dfrac{A}{d_1/\epsilon_1 + d_2/\epsilon_2} = \dfrac{\epsilon_0 \times 0.01}{5\times10^{-4} + 1.25\times10^{-4}} \\
			&= 142\ \text{pF} \\
			C_1 &= \dfrac{2\epsilon_0 A}{10^{-3}} = 177\ \text{pF} \\
			C_2 &= \dfrac{8\epsilon_0 A}{10^{-3}} = 708\ \text{pF} \\
			\dfrac{C_1 C_2}{C_1 + C_2} &= \dfrac{177\times708}{885} = 142\ \text{pF}
		}
		\end{solution}
	\end{minipage}\hfill%
	\begin{minipage}[t]{0.37\linewidth}
		\raggedleft \vspace{12pt}
		$C$ = \ansbox[1.3in]{$142$ pF}
	\end{minipage}

	\begin{minipage}[t]{0.58\linewidth}
		\part Find the field and the voltage in each layer.
		\begin{solution}[1.4in]
		\eq{
			D &= \dfrac{V}{d_1/\epsilon_1 + d_2/\epsilon_2} = \dfrac{500\,\epsilon_0}{6.25\times10^{-4}} \\
			E_1 &= \dfrac{D}{2\epsilon_0} = 400\ \text{kV/m}, \qquad V_1 = E_1 d_1 = 400\ \text{V} \\
			E_2 &= \dfrac{D}{8\epsilon_0} = 100\ \text{kV/m}, \qquad V_2 = E_2 d_2 = 100\ \text{V}
		}
		\end{solution}
	\end{minipage}\hfill%
	\begin{minipage}[t]{0.37\linewidth}
		\raggedleft \vspace{12pt}
		$E_1$ = \ansbox[1.3in]{$400$ kV/m}

		\vspace{4pt}
		$E_2$ = \ansbox[1.3in]{$100$ kV/m}
	\end{minipage}

	\part Both materials break down at $20$ MV/m. Which layer fails first as the
	voltage rises, and at what total voltage?
	\begin{solution}[1in]
	The $\epsilon_{r1} = 2$ layer, which carries four times the field of the other.
	Its field reaches $20$ MV/m when the total voltage is
	$500 \times 20\times10^{6}/4\times10^{5} = 25$ kV.
	\end{solution}

\end{parts}

\vspace{1.5em}
\noindent\textbf{Documentation:} \rule[-2pt]{0.6\linewidth}{0.4pt}

\end{questions}
